\documentclass[12pt]{article}
\usepackage[utf8]{inputenc}
\usepackage[T1]{fontenc}
\usepackage[margin=1in]{geometry}
\usepackage{setspace}
\usepackage{xcolor}
\usepackage[colorlinks=true, urlcolor=blue!60!black, linkcolor=blue!60!black]{hyperref}
\setlength{\parindent}{0em}
\setlength{\parskip}{0.45em}
\renewcommand{\baselinestretch}{1.0}
\begin{document}
\begin{center}
{\Large \textit{New York Times Co. v. Sullivan}, 376 U.S. 254 (1964)\par}
\vspace{0.3em}
{\normalsize Case Brief --- One-Page Edition\par}
\end{center}

\textbf{Memory Jogger.} A \$500,000 Alabama libel verdict over a civil-rights ad fell: public officials must prove ``actual malice'' (\href{https://supreme.justia.com/cases/federal/us/376/254/}{Justia}).

\textbf{Facts.} In 1960 the \textit{New York Times} ran ``Heed Their Rising Voices,'' a fundraising advertisement for Dr. King's defense with minor errors about Montgomery events. Commissioner L.B. Sullivan, who supervised the police but was not named, sued the \textit{Times} and four clergymen listed as endorsers.

\textbf{Procedural History.} Told the ad was libelous per se, with malice and damages presumed, the jury awarded \$500,000; the Supreme Court of Alabama affirmed.

\textbf{Votes.} Reversed, 9--0: Brennan for the Court; Black and Goldberg (each with Douglas) concurred.

\textbf{Holding.} A public official must prove with convincing clarity that a falsehood about his official conduct was made knowingly or with reckless disregard of the truth; the evidence showed neither actual malice nor that the ad concerned Sullivan.

\textbf{Reasoning.} State-court libel judgments are state action, and libel has no ``talismanic immunity.'' Errors need ``breathing space''; the Sedition Act controversy reveals the First Amendment's central meaning, and huge civil awards chill criticism. Failing to check the paper's own files was negligence at most, and imputing criticism of the police to their supervisor revived seditious libel. Black and Goldberg would make criticism of official conduct absolutely privileged. \emph{Later:} \textit{Curtis Publishing} and \textit{Gertz}.

\textbf{Examples.} \emph{Same-Side:} a careless op-ed about an Attorney General. \emph{Opposite-Side:} a blogger who purposefully avoids the truth. \emph{Fence-Sitter:} a one-time hearing witness accused of bribery.

\textbf{Critique.} Originalists say the Amendments never displaced state libel law; pragmatists say the rule under-deters lies online.

\textbf{Key Quotation.} ``The constitutional guarantees require, we think, a federal rule that prohibits a public official from recovering damages for a defamatory falsehood relating to his official conduct unless he proves that the statement was made with `actual malice'---that is, with knowledge that it was false or with reckless disregard of whether it was false or not.'' --- Justice Brennan, 376 U.S. at 279--80.

\end{document}
